%% kapa: exam=1 round=32 year=2021 period=1 form=A subject=economics subject_ko=경제학원론 seq=25 number=65
%% source: Q-Net 감정평가사 기출문제 (cst00309) artlSeq=5208397 file=extracted/1008/5208397/A형/2021년도 제32회 감정평가사 1차 1교시 A형.pdf
%% conversion: kordoc 4.24.51 -> tools/gen1.py (auto); review=proofread-sonnet; flags=none
\begin{question}{65}
화폐의 중립성이 성립하면 발생하는 현상으로 옳은 것은?
\begin{choices}
  \item 장기적으로는 고전적 이분법을 적용할 수 없다.
  \item 통화정책은 장기적으로 실업률에 영향을 줄 수 없다.
  \item 통화정책은 장기적으로 실질 경제성장률을 제고할 수 있다.
  \item 통화정책으로는 물가지수를 관리할 수 없다.
  \item 중앙은행은 국채 매입을 통해 실질 이자율을 낮출 수 있다.
\end{choices}
\end{question}
\answer{2}
